\scanpage{90}
\mspage{7}
% archive sheet number: 89
monodromie, dans les cas \deleted{\unsure{qui seront}} traités plus bas, en prouvant que l'action de $I$ sur chacun des ${\unsure{\Phi|\varphi}}^i$ est \deleted{quasi}\margin{\unsure{non!}} unipotente, \deleted{et}\inserted{en prouvant} pour ceci, que \deleted{suffira de} l'action sur chacun des \uline{faisceaux} $\ell$-adiques $\underline\Phi^i$ est \deleted{\unsure{q}} \uline{essentiellement triviale}, et mieux, qu'il existe un sous-groupe \deleted{\unsure{d'ind}} ouvert fixe $I'$ de $I$ qui opère trivialement sur \uline{tous} les \deleted{$\unsure{\Phi}$} $\underline\Phi^i_\nu$ ($i,\nu$ variables).

\deleted{Les calculs \unsure{précédents}}. C'est là un résultat de nature \uline{purement local} sur $X$, qui est accessible au calcul sous des conditions favorables, \unsure{examinées} au n° suivant.

1.21. Généralisation des développements précédents au cas d'un ouvert $U$ de la fibre générique. \inserted{($X$ tjs propre sur $S$). Soit $Y=X-U\supset X_0$.} On trouve une suite exacte des cycles évanescents
\[
\cdots\longrightarrow\Phi^i_\nu\longrightarrow H^i(X_s,\Lambda_\nu)\longrightarrow H^i(U,\Lambda_\nu)\longrightarrow\Phi^{i+1}_\nu\longrightarrow\cdots
\tag*{(1.22)}
\]
avec
\[
\Phi^i_\nu=\varinjlim_\alpha H^i_{Y_\alpha}(\deleted{\unsure{X}}X_\alpha,\Lambda_\nu),
\tag*{(1.23)}
\]
où $Y_\alpha$ est l'image inverse de $Y_\alpha$ dans $X_\alpha$. La suite spectrale de localisation est
\[
\Phi^*_\nu\Longleftarrow E_2^{pq}=H^p(X_0,\underline\Phi^q_\nu),
\tag*{(1.24).}
\]
\[
\underline\Phi^q_\nu=\varinjlim_\alpha\bigl(\underline H^q_{Y_\alpha}((\Lambda_\nu)_{X_\alpha})\mid X_{\alpha0}\bigr).
\tag*{(1.25)}
\]
